\addcontentsline{toc}{subsubsection}{Solution 13. The Fixed-Point Lever}

\begin{dlsummary}[title={\bfseries Solution 13. The Fixed-Point Lever},breakable]

\begin{proof}
Setting $y=x$ in the given inequality gives
\[
f(x)\bigl(x+f(x)\bigr)
\geq f(f(x))\bigl(f(x)+x\bigr).
\]
Since $x+f(x)>0$, it follows that
\[
f(f(x))\leq f(x). \tag{1}
\]

\medskip
\noindent
\textbf{Step 1. Use the assumed two-cycle.}

Let $b=f(a)$. Then $f(b)=a$. Applying (1) at $a$ gives
$a\leq b$, while applying it at $b$ gives
\[
b=f(f(b))\leq f(b)=a.
\]
Thus $a=b$ and $f(a)=a$.

Now set $x=a$ in the original inequality. Since
$f(a)=f(f(a))=a$, we get
\[
a\bigl(a+f(y)\bigr)\geq a(a+y),
\]
so
\[
f(y)\geq y\qquad(y>0). \tag{2}
\]

\medskip
\noindent
\textbf{Step 2. Force idempotence.}

Applying (2) to $f(x)>0$ yields $f(f(x))\geq f(x)$.
Together with (1), this gives
\[
f(f(x))=f(x). \tag{3}
\]

Substituting (3) into the original inequality and dividing by
$f(x)>0$ yields
\[
x+f(y)\geq f(x)+y.
\]
Therefore $f(x)-x\leq f(y)-y$ for all $x,y>0$.
Interchanging $x$ and $y$ shows that $f(x)-x$ is a constant
$c$. By (2), $c\geq0$.

Finally, $f(x)=x+c$ and (3) imply
\[
x+2c=f(f(x))=f(x)=x+c,
\]
so $c=0$. Hence $f(x)=x$ for every $x>0$.

\medskip
\noindent
\textbf{Verification.}

For $f(x)=x$, the inequality holds with equality, and
$f(f(a))=a$ for every $a>0$.
\end{proof}

\end{dlsummary}